\documentclass[12pt]{article}
\usepackage[utf8]{inputenc}
\usepackage[T1]{fontenc}
\usepackage{hyperref}
\title{Yakaikeke Language Resources for K-12 Curriculum Development}
\author{Compiled by Claude Sonnet 5.5 (claude-sonnet-5-5) for Andr\'as Kornai}
\date{}
\begin{document}
\maketitle

\section{Introduction}
Yakaikeke (ISO 639-3 ykk; also spelled Iakaikeke) is an endangered Oceanic language of Milne Bay Province, Papua New Guinea. SIL (1998) gave about 100 native speakers; Joshua Project lists about 200 people. Very little material exists, and none of it is educational.

\section{Primary Educational Resources}
None found.

\section{Linguistic and Reference Materials}
\begin{itemize}
\item Yakaikeke language (Wikipedia): \url{https://en.wikipedia.org/wiki/Yakaikeke_language}
\item Numeral systems of Yakaikeke (MPI-EVA): \url{https://lingweb.eva.mpg.de/channumerals/Yakaikeke.htm}
\item Endangered Languages Project entry: \url{https://endangeredlanguages.com/lang/3188?hl=es}
\end{itemize}

\section{Religious and Community Resources}
\begin{itemize}
\item Joshua Project, Yakaikeke people group (no Bible translation begun): \url{https://joshuaproject.net/people_groups/15954}
\end{itemize}

\section{Digital and Audio Resources}
\begin{itemize}
\item Global Recordings Network language page for ykk (contents were not checked): \url{https://globalrecordings.net/en/language/ykk}
\end{itemize}

\section{Curriculum Development Context}
A 2009 note by a SIL linguist reports that speakers have largely shifted to English numerals. Joshua Project suggests the group may be assimilating into larger neighbours. A native-language curriculum would need documentation first.
\end{document}